\section{Supplementary results and proofs for \texorpdfstring{Sections~\ref{sec:sensitive}--\ref{sec:smallarea}}{Sections 5\textendash 6}}\label{app:proofs-design}

\providecommand{\armT}{\mathsf{T}}
\providecommand{\armC}{\mathsf{C}}
\providecommand{\armD}{\mathsf{D}}
\providecommand{\armL}{\mathsf{L}}
\providecommand{\SDB}{\mathrm{SDB}}
\providecommand{\ECV}{\mathrm{ECV}}

This appendix states the supplementary results to which \cref{sec:sensitive,sec:smallarea} point and proves every result of those sections. Throughout, $\mathcal A=(A_1,\dots,A_n)$ is the vector of arm assignments. Under \cref{asm:le}(i), the vectors $(W_i,Z_i,D_i)$, $i=1,\dots,n$, are independent and identically distributed and independent of $\mathcal A$. Conditionally on $\mathcal A$, the respondents in each arm therefore form an independent random sample from the population, and samples in different arms are independent. We use this fact without further comment.

\subsection{The variance of the list experiment}\label{app:pd-listvar}

\begin{remark}[Finite-population variance]\label{rem:le-neyman}
Condition on the $n$ sampled respondents and take the in-sample prevalence $\pi_n=n^{-1}\sum_iZ_i$ as the estimand. Under complete randomisation, Neyman's formula \citep[ch.~6]{imbens2015} gives
$\Var(\widehat\pi\mid\text{sample})=S_{\armT}^2/n_{\armT}+S_{\armC}^2/n_{\armC}-\pi_n(1-\pi_n)/(n-1)$,
where $S_{\armT}^2$ and $S_{\armC}^2$ are the in-sample variances (divisor $n-1$) of $W_i+Z_i$ and $W_i$. In a general experiment the last term, the variance of unit-level effects, is not identified, and the usual variance estimator is conservative. In a list experiment the unit-level effect is $Z_i\in\{0,1\}$, so the term is a known function of $\pi_n$ and can be estimated by $\widehat\pi(1-\widehat\pi)/(n-1)$.
\end{remark}

\emph{Proof of \cref{prop:list-var}.} \emph{Unbiasedness and \eqref{eq:le-var}.} Under \cref{asm:le}(ii)--(iii), $\overline Y_{\armT}$ is the mean of $n_{\armT}$ independent copies of $W_i+Z_i$ and $\overline Y_{\armC}$ the mean of $n_{\armC}$ independent copies of $W_i$, and the two means are independent. Hence $\E\widehat\pi=\E(W_i+Z_i)-\E W_i=\pi$ and $\Var(\widehat\pi)=\Var(W_i+Z_i)/n_{\armT}+\Var(W_i)/n_{\armC}$. Expanding $\Var(\sum_jW_{ij})$ and $\Var(\sum_jW_{ij}+Z_i)$ by bilinearity gives \eqref{eq:le-sig}.

\emph{Part (a).} Suppose first $\sigma_{\armT},\sigma_{\armC}>0$. The function $r\mapsto\sigma_{\armT}^2/r+\sigma_{\armC}^2/(1-r)$ is strictly convex on $(0,1)$, and its derivative $-\sigma_{\armT}^2/r^2+\sigma_{\armC}^2/(1-r)^2$ vanishes only at $r/(1-r)=\sigma_{\armT}/\sigma_{\armC}$, that is at $r^*$. There its value is $\sigma_{\armT}(\sigma_{\armT}+\sigma_{\armC})+\sigma_{\armC}(\sigma_{\armT}+\sigma_{\armC})=(\sigma_{\armT}+\sigma_{\armC})^2$, which gives \eqref{eq:le-loss}. If one of the two variances is zero, \eqref{eq:le-loss} is the infimum of $L$, approached as $r$ tends to $0$ or $1$.

\emph{Part (b).} Write $U=(W_i+Z_i)-\E(W_i+Z_i)$ and $V=-(W_i-\E W_i)$, so that $U+V=Z_i-\pi$. By Minkowski's inequality in $L^2$,
\[
  \sqrt{\pi(1-\pi)}=\|U+V\|_2\le\|U\|_2+\|V\|_2=\sigma_{\armT}+\sigma_{\armC},
\]
so $L^*\ge1$. Equality holds if and only if $V=0$ almost surely, or $U=cV$ almost surely for some $c\ge0$. The first case is $W_i$ constant. In the second, $Z_i-\pi=-(1+c)(W_i-\E W_i)$, so $W_i=\E W_i-(Z_i-\pi)/(1+c)$. Because $0<\pi<1$, $W_i$ then takes two values that differ by $1/(1+c)\in(0,1]$. Since $W_i$ is integer-valued, $1/(1+c)=1$, so $c=0$ and $W_i+Z_i$ is constant. Conversely, if $W_i\equiv w_0$ then $\sigma_{\armC}=0$ and $\sigma_{\armT}^2=\pi(1-\pi)$; if $W_i+Z_i\equiv y_0$ then $\sigma_{\armT}=0$ and $\sigma_{\armC}^2=\pi(1-\pi)$. In both cases $L^*=1$. In the first case the treatment count $w_0+Z_i$ reveals $Z_i$, and in the second the control count $y_0-Z_i$ does. The converse fails. If $W_i$ is uniform on $\{0,2\}$ and independent of $Z_i$, the parity of $W_i+Z_i$ reveals $Z_i$, but $\sigma_{\armC}^2=1$ and $\sigma_{\armT}^2=1+\pi(1-\pi)$, so $L^*>1$; for $\pi=0.2$, $L^*\approx27$.

\emph{Part (c).} Each $c_{jk}$ (which appears twice in the sum over $j\ne k$) and each $\gamma_j$ enters $\sigma_{\armC}^2$ or $\sigma_{\armT}^2$ linearly with a positive coefficient, and neither enters with a negative one. Since $x\mapsto\sqrt x$ is strictly increasing, $\sigma_{\armT}+\sigma_{\armC}$ and hence $L^*$ are strictly increasing in each. If $\sum_j\gamma_j\ge0$ then $\sigma_{\armT}\ge\sigma_{\armC}$, and
\[
  (\sigma_{\armT}+\sigma_{\armC})^2\ge\sigma_{\armT}^2+3\sigma_{\armC}^2=4\sigma_{\armC}^2+\pi(1-\pi)+2\textstyle\sum_j\gamma_j\ge4\sigma_{\armC}^2+\pi(1-\pi).
\]
Dividing by $\pi(1-\pi)$ gives the bound. For exchangeable items, $\sigma_{\armC}^2=Jp(1-p)+J(J-1)\rho p(1-p)$, which is \eqref{eq:le-exch}. The condition $\rho\ge-1/(J-1)$ is $\sigma_{\armC}^2\ge0$, which is necessary. It is not sufficient for binary items. For $p=1/2$ and odd $J$, the integer-valued $W_i$ has mean $J/2$, so $\sigma_{\armC}^2\ge1/4$, which gives $\rho\ge-1/J$. This bound is attained by drawing $W_i$ uniformly from $\{(J-1)/2,(J+1)/2\}$ and endorsing a uniformly random subset of that size. For large $J$, $(\sigma_{\armT}+\sigma_{\armC})^2=4\sigma_{\armC}^2+O(1)$ when $\sum_j\gamma_j$ is bounded, which gives the growth rate stated for uncorrelated items.

\emph{\Cref{rem:le-neyman}.} Condition on the sample, so that the $(W_i,Z_i)$ are fixed and only $\mathcal A$ is random. Under \cref{asm:le}(ii)--(iii) the potential responses are $Y_i(\armC)=W_i$ and $Y_i(\armT)=W_i+Z_i$, and the unit-level effect is $Y_i(\armT)-Y_i(\armC)=Z_i$. Neyman's formula for completely randomised experiments \citep[ch.~6]{imbens2015} gives $\Var(\widehat\pi\mid\text{sample})=S_{\armT}^2/n_{\armT}+S_{\armC}^2/n_{\armC}-S_Z^2/n$, where $S_Z^2=(n-1)^{-1}\sum_i(Z_i-\pi_n)^2$. Because $Z_i^2=Z_i$, $\sum_i(Z_i-\pi_n)^2=n\pi_n(1-\pi_n)$, so $S_Z^2/n=\pi_n(1-\pi_n)/(n-1)$. \qed

\subsection{Proof of \texorpdfstring{\cref{thm:joint-sdb}}{the joint SDB theorem}}\label{app:pd-joint}

We first record the identity $\SDB=\pi^H-\pi^F$:
\[
  \pi-\pi^D=\{\Pr(Z_i=1,D_i=0)+\Pr(Z_i=1,D_i=1)\}-\{\Pr(Z_i=0,D_i=1)+\Pr(Z_i=1,D_i=1)\}=\pi^H-\pi^F.
\]

\begin{lemma}[Stratified means in the paired design]\label{lem:pd-strat}
Under \cref{asm:le}(i)--(iv), suppose $n_a/n\to r_a>0$ for $a\in\{\armC,\armT\}$ and $0<P_0<1$. For $a\in\{\armC,\armT\}$ and $d\in\{0,1\}$, let $m_{ad}=\E(Y_i\mid a,D_i=d)$ and $v_{ad}=\Var(Y_i\mid a,D_i=d)$, so $m_{\armT d}-m_{\armC d}=\Delta_d$, $v_{\armT d}=\Var(W_i+Z_i\mid D_i=d)$ and $v_{\armC d}=\Var(W_i\mid D_i=d)$. Then $\sqrt n(\widehat P_0-P_0)$, $\sqrt n(\overline Y_{ad}-m_{ad})$ for the four pairs $(a,d)$ are jointly asymptotically normal and asymptotically independent, with asymptotic variances $P_0P_1$ and $v_{ad}/(r_aP_d)$.
\end{lemma}

\begin{proof}
Within arm $a$, $\overline Y_{ad}=\{n_a^{-1}\sum_{i\in a}\ind{D_i=d}Y_i\}/\{n_a^{-1}\sum_{i\in a}\ind{D_i=d}\}$ is a ratio of two sample means of independent and identically distributed terms. By the multivariate central limit theorem and the delta method, $\sqrt{n_a}(\overline Y_{ad}-m_{ad})=n_a^{-1/2}\sum_{i\in a}\ind{D_i=d}(Y_i-m_{ad})/P_d+o_p(1)$. The summands for $d=0$ and $d=1$ have zero covariance, because $\ind{D_i=0}\ind{D_i=1}=0$. Each also has zero covariance with $\ind{D_i=0}-P_0$, since $\E[\ind{D_i=d}(Y_i-m_{ad})(\ind{D_i=0}-P_0)]=(\ind{d=0}-P_0)\E[\ind{D_i=d}(Y_i-m_{ad})]=0$. The estimator $\widehat P_0$ pools all arms, which have the same law of $D_i$ by \cref{asm:le}(iv), and $\sqrt n(\widehat P_0-P_0)=n^{-1/2}\sum_{i=1}^n(\ind{D_i=0}-P_0)$. Arms are independent. The joint central limit theorem for the stacked influence terms, with $n_a=r_an+o(n)$, gives the claim, and the variance of $\ind{D_i=d}(Y_i-m_{ad})/P_d$ is $v_{ad}/P_d$.
\end{proof}

\emph{Part (a).} The three arm means are independent sample means with expectations $\E(W_i+Z_i)$, $\E W_i$ and $\pi^D$ and variances $\sigma_{\armT}^2/n_{\armT}$, $\sigma_{\armC}^2/n_{\armC}$ and $\pi^D(1-\pi^D)/n_{\armD}$, which proves unbiasedness and \eqref{eq:le-split}; $\pi$ and $\pi^D$ are identified as the limits of these means. In the split design, the law of the observed data depends on the joint law $Q$ of $(W_i,Z_i,D_i)$ only through the law of $(W_i,Z_i)$ (list arms) and the law of $D_i$ (direct arm). Let $t=\Pr(Z_i=1,D_i=1)$. The Fréchet--Hoeffding bounds for a $2\times2$ table with margins $\pi$ and $\pi^D$ give $\max(0,\pi+\pi^D-1)\le t\le\min(\pi,\pi^D)$, and $\pi^H=\pi-t$, which yields \eqref{eq:le-bounds}. For sharpness, fix any $t$ in this range. Define $Q'$ by taking $(Z_i,D_i)$ with $\Pr(Z_i=1,D_i=1)=t$ and margins $\pi,\pi^D$, which is a valid table because $t$ lies in the Fréchet range, and then drawing $W_i$ from its $Q$-conditional law given $Z_i$, independently of $D_i$. Then $Q'$ and $Q$ induce the same law of $(W_i,Z_i)$ and of $D_i$, hence the same law of the observed data, and $\pi^H(Q')=\pi-t$. Every point of \eqref{eq:le-bounds} is therefore compatible with the data. Under \cref{asm:le-mono}, $\Pr(Z_i=0,D_i=1)=0$ forces $t=\pi^D$, and so $\pi^H=\pi-\pi^D=\SDB$.

\emph{Part (b).} By \cref{asm:le}(iv), $D_i$ is the same variable in every arm, and by (i) the assignment is independent of $(W_i,Z_i,D_i)$. Hence $\E(Y_i\mid\armT,D_i=d)=\E(W_i+Z_i\mid D_i=d)$ and $\E(Y_i\mid\armC,D_i=d)=\E(W_i\mid D_i=d)$, whose difference is $\E(Z_i\mid D_i=d)=\Delta_d$; this is \eqref{eq:le-strat}. The formulas $\pi^H=P_0\Delta_0$, $\pi^F=P_1(1-\Delta_1)$ and $\pi=\sum_dP_d\Delta_d$ follow from the definitions, and $P_d$ is identified from $D_i$ in any arm. Consistency of $\widehat{\SDB}_{\mathrm{PS}}=g(\widehat P_0,\widehat\Delta_0,\widehat\Delta_1)$, with $g(P_0,\Delta_0,\Delta_1)=P_0\Delta_0+(1-P_0)\Delta_1-(1-P_0)$ and $\widehat\Delta_d=\overline Y_{\armT d}-\overline Y_{\armC d}$, follows from \cref{lem:pd-strat}. The partial derivatives of $g$ are $\partial g/\partial P_0=1+\Delta_0-\Delta_1$ and $\partial g/\partial\Delta_d=P_d$. By \cref{lem:pd-strat}, $\sqrt n(\widehat\Delta_d-\Delta_d)$ has asymptotic variance $(v_{\armT d}/r_{\armT}+v_{\armC d}/r_{\armC})/P_d$, independently across $d$ and of $\widehat P_0$. The delta method gives \eqref{eq:le-vps}.

\emph{Part (c).} Under \cref{asm:le-mono}, $D_i=1$ implies $Z_i=1$, so $\Delta_1=1$, $\pi^F=0$ and $\SDB=\pi^H=P_0\Delta_0$.

(i) Let $\mathcal G$ collect $\mathcal A$ and all direct answers, and let $E$ be the event that $n_{\armT0}\ge1$ and $n_{\armC0}\ge1$, where $n_{a0}$ counts respondents in arm $a$ with $D_i=0$. Conditionally on $\mathcal G$, the list answers of respondents in arm $a$ with $D_i=0$ are independent draws from the law of $Y_i$ given $(a,D_i=0)$. Hence on $E$, $\E(\widehat\pi^H\mid\mathcal G)=\widehat P_0\Delta_0$, while $\widehat\pi^H=0$ on $E^c$. Therefore $\E\widehat\pi^H=\Delta_0\E(\widehat P_0\ind{E})=\pi^H-\Delta_0\E(\widehat P_0\ind{E^c})$, and
\[
  |\E\widehat\pi^H-\pi^H|\le\Pr(E^c)\le\Pr(n_{\armT0}=0)+\Pr(n_{\armC0}=0)=P_1^{\,n_{\armT}}+P_1^{\,n_{\armC}}.
\]

(ii) Apply the delta method to $(P_0,\Delta_0)\mapsto P_0\Delta_0$ with \cref{lem:pd-strat}: the asymptotic variance is $\Delta_0^2P_0P_1+P_0^2(v_{\armT0}/r_{\armT}+v_{\armC0}/r_{\armC})/P_0$, which is \eqref{eq:le-vh}.

(iii) Under \cref{asm:le,asm:le-mono} the parameter is the joint law $Q$ of $(W_i,Z_i,D_i)$ on $\mathcal S=\{0,\dots,J\}\times\{0,1\}^2\setminus\{(w,0,1):w\le J\}$. Reparametrise $Q$ by $p_0=Q(D_i=0)$, $f(w,z)=Q(W_i=w,Z_i=z\mid D_i=0)$ and $g(w)=Q(W_i=w\mid D_i=1)$; under monotonicity $Z_i=1$ on $\{D_i=1\}$, so this is a bijection. Let $q_{\armC}(y)=f(y,0)+f(y,1)$ and $q_{\armT}(y)=f(y,0)+f(y-1,1)$ be the laws of the list count in the two arms within $\{D_i=0\}$, with cumulative distribution functions $F_{\armC}$ and $F_{\armT}$. By the computation that gives \eqref{eq:le-types}, $f(y,1)=F_{\armC}(y)-F_{\armT}(y)$ and $f(y,0)=F_{\armT}(y)-F_{\armC}(y-1)$. Conversely, any pair $(q_{\armC},q_{\armT})$ of laws on $\{0,\dots,J\}$ and $\{0,\dots,J+1\}$ with $F_{\armC}(y-1)\le F_{\armT}(y)\le F_{\armC}(y)$ for all $y$ defines through these formulas a law $f$ whose induced pair is $(q_{\armC},q_{\armT})$: the two formulas sum to $q_{\armC}(y)$, and $f(y,0)+f(y-1,1)=q_{\armT}(y)$. So $f\leftrightarrow(q_{\armC},q_{\armT})$ is a bijection onto the set $\mathcal Q$ cut out by these inequalities, and interior points $f>0$ correspond to strict inequalities. The likelihood of the observed data is
\[
  p_0^{N_0}(1-p_0)^{N_1}\prod_yq_{\armC}(y)^{n_{\armC0}(y)}\prod_yq_{\armT}(y)^{n_{\armT0}(y)}\prod_yg(y)^{n_{\armC1}(y)}g(y-1)^{n_{\armT1}(y)},
\]
where $N_d$ counts all respondents with $D_i=d$ and $n_{ad}(y)$ counts those in arm $a$ with $D_i=d$ and $Y_i=y$. It factorises over $p_0$, $(q_{\armC},q_{\armT})\in\mathcal Q$ and $g$, which vary independently. Without the constraint $\mathcal Q$, the maximisers of the first three factors are the empirical frequencies. At an interior point these converge almost surely to a point of the interior of $\mathcal Q$, so with probability tending to one they satisfy the constraint and are the constrained maximisers. Since $\pi^H=p_0\sum_yy\{q_{\armT}(y)-q_{\armC}(y)\}$, invariance of maximum likelihood gives $\widehat\pi^H_{\mathrm{ML}}=\widehat P_0(\overline Y_{\armT0}-\overline Y_{\armC0})=\widehat\pi^H$ on that event. Locally the model is a smooth parametric family with an open parameter set and non-singular Fisher information (a product of multinomials with positive cell probabilities). The arm samples are independent and identically distributed with sizes proportional to $n$, so the experiment is locally asymptotically normal. The maximum likelihood estimator of the smooth functional $\pi^H$ then attains the Cramér--Rao bound, and by the convolution theorem no regular estimator has smaller asymptotic variance \citep[chs.~7--8]{vandervaart1998}. Its asymptotic variance is $V_H$ by (ii).

(iv) Under \cref{asm:le-mono}, $\Delta_1=1$, so the first term of \eqref{eq:le-vps} equals the first term of \eqref{eq:le-vh}. The difference is $P_1(v_{\armT1}/r_{\armT}+v_{\armC1}/r_{\armC})$, and $v_{\armT1}=\Var(W_i+1\mid D_i=1)=v_{\armC1}$.

\emph{Part (d).} Without \cref{asm:le-mono}, \cref{lem:pd-strat} still gives $\widehat P_0\to P_0$ and $\widehat\Delta_0\to\Delta_0$, so $\widehat\pi^H\to P_0\Delta_0=\pi^H=\SDB+\pi^F$. Also $\widehat P_1+\widehat\pi^H\to\Pr(D_i=1)+\Pr(Z_i=1,D_i=0)=\pi+\pi^F$, since $\Pr(D_i=1)=\Pr(Z_i=1,D_i=1)+\pi^F$. Consistency of $\widehat{\SDB}_{\mathrm{PS}}$ was shown in (b), and the split-design statement is part (a). \qed

\subsection{Testable implications and covariates}\label{app:pd-tests}

\begin{corollary}[Testable implications]\label{cor:le-tests}
\begin{enumerate}[(a)]
  \item \emph{Design effects} \citep{blair2012}. Under \cref{asm:le}(i)--(iii), for $y=0,\dots,J$,
  \begin{equation}\label{eq:le-types}
  \begin{aligned}
    \Pr(W_i=y,Z_i=1)&=\Pr(Y_i\le y\mid\armC)-\Pr(Y_i\le y\mid\armT),\\
    \Pr(W_i=y,Z_i=0)&=\Pr(Y_i\le y\mid\armT)-\Pr(Y_i\le y-1\mid\armC),
  \end{aligned}
  \end{equation}
  and both must be non-negative. In the paired design under \cref{asm:le}(iv), the same holds within each stratum $D_i=d$.
  \item \emph{Monotonicity.} Under \cref{asm:le} and \cref{asm:le-mono} in the paired design, $\Pr(Y_i=y+1\mid\armT,D_i=1)=\Pr(Y_i=y\mid\armC,D_i=1)$ for every $y$. In particular $\Delta_1=1$, which is the first placebo test of \citet{aronow2015}.
  \item \emph{Treatment independence.} Under \cref{asm:le}(iv), the law of $D_i$ is the same in $\armD$, $\armC$ and $\armT$. Comparing $\armT$ with $\armC$ is the second placebo test of \citet{aronow2015}; the direct arm adds a comparison with respondents who never saw the list.
\end{enumerate}
\end{corollary}

The test of \citet{blair2012} estimates the $2(J+1)$ proportions in \eqref{eq:le-types} and rejects if any is significantly negative, with a Bonferroni correction. Like every test in the corollary, it detects only violations that change the law of the observed data, and design effects of opposite signs can cancel. Passing these tests is necessary for the assumptions, not sufficient.

\emph{Proof of \cref{cor:le-tests}.} (a) Under \cref{asm:le}(i)--(iii), $Y_i$ has the law of $W_i$ in $\armC$ and of $W_i+Z_i$ in $\armT$. The event $\{W_i+Z_i\le y\}$ is the disjoint union of $\{W_i\le y-1\}$ and $\{W_i=y,Z_i=0\}$, so $\Pr(Y_i\le y\mid\armT)=\Pr(Y_i\le y-1\mid\armC)+\Pr(W_i=y,Z_i=0)$, which is the second identity in \eqref{eq:le-types}. Subtracting it from $\Pr(W_i=y)=\Pr(Y_i\le y\mid\armC)-\Pr(Y_i\le y-1\mid\armC)$ gives the first. Under \cref{asm:le}(iv) the assignment is independent of $(W_i,Z_i,D_i)$, so the argument applies conditionally on $D_i=d$.
(b) Under \cref{asm:le-mono}, $Z_i=1$ on $\{D_i=1\}$, so within that stratum $Y_i=W_i+1$ in $\armT$ and $Y_i=W_i$ in $\armC$, and $W_i$ has the same conditional law in both arms.
(c) Under \cref{asm:le}(i) and (iv), $D_i$ is a single variable independent of the assignment. \qed

\begin{remark}[Covariates and a parametric joint model]\label{rem:le-joint}
All statements in \cref{thm:joint-sdb} hold conditionally on discrete pre-treatment covariates $X_i=x$, so they identify $\SDB(x)$, $\pi^H(x)$ and $\pi^F(x)$. A parametric joint model in the spirit of \citet{imai2011} writes, under \cref{asm:le-mono},
\begin{equation}\label{eq:le-joint}
\begin{gathered}
  \Pr(Z_i=1\mid x)=\logit^{-1}(x^\top\beta),\qquad
  \Pr(D_i=1\mid Z_i=1,x)=\logit^{-1}(x^\top\zeta),\\
  W_i\mid(Z_i=z,X_i=x)\sim q_z(\cdot\,;x,\psi),
\end{gathered}
\end{equation}
where $q_z$ is a binomial or beta-binomial count model \citep{imai2011} and $\alpha(x)=\logit^{-1}(x^\top\zeta)$ is the \emph{candour rate}, the probability that a private endorser endorses directly. Then $\SDB(x)=\pi(x)\{1-\alpha(x)\}$. In the split design the direct arm identifies $\pi(x)\alpha(x)$ and the list arms identify $\pi(x)$, so $\beta$ and $\zeta$ are identified when $X_i$ contains an intercept and its support spans $\mathbb R^{\dim X_i}$. In the paired design the likelihood of $(D_i,Y_i)$ is a finite mixture over $Z_i$ with known structure. If social pressure on an item runs towards agreement for some respondents, \cref{asm:le-mono} is replaced by its mirror image, $\Pr(Z_i=1,D_i=0)=0$; the theorem then holds with the roles of $D_i=0$ and $D_i=1$ exchanged, and $\SDB=-\pi^F\le0$.
\end{remark}

\subsection{Optimal design}\label{app:pd-design}

A direct interview costs $c_{\armD}$ and an interview in a list arm costs $c_{\armL}\ge c_{\armD}$; in the paired design $c_{\armL}$ includes the direct question. The budget is $B=c_{\armD}n_{\armD}+c_{\armL}(n_{\armC}+n_{\armT})$, and arm sizes are treated as continuous. For the paired design under \cref{asm:le-mono} we write the variance \eqref{eq:le-vh} in absolute terms as $\nu_0/n+b_{\armT}/n_{\armT}+b_{\armC}/n_{\armC}$, with
\[
  \nu_0=\Delta_0^2P_0P_1,\qquad b_{\armT}=P_0v_{\armT0},\qquad b_{\armC}=P_0v_{\armC0},\qquad \bar b=\sqrt{b_{\armT}}+\sqrt{b_{\armC}}.
\]
\emph{Floor and ceiling effects} \citep{blair2012} are modelled as follows: in arm $\armT$, a respondent with $Z_i=1$ and $W_i=J$ reports $J$ instead of $J+1$ with probability $\phi_{\mathrm c}$, and a respondent with $Z_i=1$ and $W_i=0$ reports $0$ instead of $1$ with probability $\phi_{\mathrm f}$. For a candidate control list $\ell$ (its statements and hence $J$), let $V^*_\ell$ be $B$ times the smallest variance attainable with $\ell$, and let $\operatorname{bias}_\ell$ be the resulting bias of the SDB estimator. The \emph{exposure} of a list is $e_\ell=\Pr(W_i\in\{0,J\})$.

\begin{theorem}[Optimal allocation and choice of the control list]\label{thm:design}
\begin{enumerate}[(a)]
  \item \emph{Split design.} With $v_{\armD}=\pi^D(1-\pi^D)$, $v_{\armC}=\sigma_{\armC}^2$ and $v_{\armT}=\sigma_{\armT}^2$, the variance \eqref{eq:le-split} is minimised subject to the budget by
  \begin{equation}\label{eq:le-neyman}
    n_a=\frac{B}{K_{\mathrm S}}\sqrt{\frac{v_a}{c_a}}\quad(a\in\{\armD,\armC,\armT\},\ c_{\armC}=c_{\armT}=c_{\armL}),\qquad K_{\mathrm S}=\sum_a\sqrt{v_ac_a},\qquad V^*=K_{\mathrm S}^2 .
  \end{equation}
  With equal costs the arm shares are proportional to $(\sqrt{v_{\armD}},\sigma_{\armC},\sigma_{\armT})$.
  \item \emph{Paired design.} Let $c_+=c_{\armL}-c_{\armD}$, the marginal cost of moving a respondent from the direct arm into a list arm. If $\nu_0c_+>c_{\armD}\bar b^2$, the optimum uses all three arms:
  \begin{equation}\label{eq:le-paired-int}
    n=\frac{B}{K_{\mathrm P}}\sqrt{\frac{\nu_0}{c_{\armD}}},\quad
    n_{\armT}=\frac{B}{K_{\mathrm P}}\sqrt{\frac{b_{\armT}}{c_+}},\quad
    n_{\armC}=\frac{B}{K_{\mathrm P}}\sqrt{\frac{b_{\armC}}{c_+}},\quad
    K_{\mathrm P}=\sqrt{\nu_0c_{\armD}}+\bar b\sqrt{c_+},\quad V^*=K_{\mathrm P}^2,
  \end{equation}
  with $n_{\armD}=n-n_{\armT}-n_{\armC}>0$. Otherwise the direct arm is empty, $n_{\armT}+n_{\armC}=B/c_{\armL}$, $n_{\armT}/n_{\armC}=\sqrt{b_{\armT}/b_{\armC}}$, and $V^*=c_{\armL}(\nu_0+\bar b^2)$.
  \item \emph{Choice of list.} Let $\mathcal L$ be a finite menu of control lists, and let $\ell(B)$ minimise $\mathrm{MSE}_\ell(B)=V^*_\ell/B+\operatorname{bias}_\ell^2$ over $\mathcal L$. Under the floor-and-ceiling model,
  \begin{equation}\label{eq:le-fcbias}
    \operatorname{bias}_\ell=-\pi\bigl\{\phi_{\mathrm c}\Pr(W_i=J\mid Z_i=1)+\phi_{\mathrm f}\Pr(W_i=0\mid Z_i=1)\bigr\}
  \end{equation}
  in the split design, and the same expression with $\pi$ replaced by $\pi^H$ and the conditioning event by $\{Z_i=1,D_i=0\}$ in the paired design with estimator \eqref{eq:le-comb}. If $B<B'$, then $V^*_{\ell(B')}\ge V^*_{\ell(B)}$ and $\operatorname{bias}_{\ell(B')}^2\le\operatorname{bias}_{\ell(B)}^2$: larger budgets select lists that are noisier but less exposed to floor and ceiling effects. If the control count is independent of $(Z_i,D_i)$, then $|\operatorname{bias}_\ell|\le\max(\phi_{\mathrm c},\phi_{\mathrm f})\,\pi\,e_\ell$ (split) or $\le\max(\phi_{\mathrm c},\phi_{\mathrm f})\,\pi^H\,e_\ell$ (paired).
  \item \emph{Structure of the control list.} Suppose the control count is independent of $(Z_i,D_i)$.
  \begin{enumerate}[(i)]
    \item For $J\ge2$ independent control items with common prevalence $p$, the exposure $p^J+(1-p)^J$ is uniquely minimised at $p=1/2$, where the variance $Jp(1-p)$ is maximised.
    \item Let $J=2m$ items form $m$ independent pairs, each item with prevalence $1/2$ and within-pair correlation $\rho\in[-1,1]$. Then
    \begin{equation}\label{eq:le-pairs}
      \sigma_{\armC}^2=\frac{m(1+\rho)}2,\qquad \Pr(W_i=0)=\Pr(W_i=J)=\Bigl(\frac{1+\rho}4\Bigr)^m,
    \end{equation}
    and both are strictly increasing in $\rho$. Hence $V^*_\ell$, $e_\ell$ and $\mathrm{MSE}_\ell(B)$ are all minimised by the most negative admissible $\rho$.
    \item At $\rho=-1$ the count is constant and the treatment count reveals $Z_i$ (\cref{prop:list-var}(b)), so $\rho$ must be restricted by a disclosure requirement. For any closed set $R\subset(-1,1]$ of admissible correlations, for example those meeting a bound on the mutual information $I(Y_i;Z_i\mid\armT)$ or on the share of treated respondents whose posterior $\Pr(Z_i=1\mid Y_i,\armT)$ exceeds a threshold, the MSE-optimal pair design uses $\rho^*=\min R$. Under the variance floor $\sigma_{\armC}^2\ge v$ with $0<v\le m$, $\rho^*=2v/m-1$, $\sigma_{\armC}^2=v$, and the exposure is $2\{v/(2m)\}^m$, which decreases in $m$.
  \end{enumerate}
\end{enumerate}
\end{theorem}

Part (a) is a Neyman-type allocation. For the inputs of the illustration in \cref{sec:design-list} ($\sigma_{\armC}^2=\sigma_{\armT}^2=1$, $\pi^D=0.2$, equal costs), equal allocation gives $n\Var=6.48$ against $5.76$ at the optimum, with shares $(1/6,5/12,5/12)$. Equal count variances implicitly assume that the control items covary slightly negatively with the sensitive item ($2\sum_j\gamma_j=-\pi(1-\pi)$). With the independent items of \cref{sec:sens-var} and $\pi=0.2$, $\sigma_{\armT}^2=1.16$; the shares become $(0.16,0.40,0.43)$, and $n\Var$ is $6.14$ against $6.96$ under equal allocation.

Part (b) has a sharper message. In the paired design every list respondent also answers the direct question, so the direct-only arm only refines $\widehat P_0$. Its contribution $\nu_0=(\pi^H)^2P_1/P_0$ is second order in $\pi^H$. The direct-only arm is used if and only if $c_{\armD}/(c_{\armL}-c_{\armD})<\nu_0/\bar b^2$. By Minkowski's inequality $\sqrt{v_{\armT0}}+\sqrt{v_{\armC0}}\ge\{\Delta_0(1-\Delta_0)\}^{1/2}$, so for every control list
\[
  \frac{\nu_0}{\bar b^2}\le\frac{\pi^HP_1}{P_0-\pi^H},
\]
and if $v_{\armC0}\ge v>0$ the threshold is at most $\Delta_0^2P_1/v=O\{(\pi^H)^2\}$, or $\Delta_0^2P_1/(4v)$ if also $v_{\armT0}\ge v$. With $\pi^H=0.05$, $\pi=0.2$, $\sigma_{\armC}^2=1$ and a control count independent of $(Z_i,D_i)$, the universal bound is $0.0094$ and the exact threshold is $1.3\times10^{-4}$. A direct-only interview would have to cost about $10^{-4}$ times the extra cost of the list module for the arm to pay. A small direct arm may still be fielded to test treatment independence (\cref{cor:le-tests}(c)), but not for precision.

Part (d) needs a caveat. Neither the variance floor nor the exposure $e_\ell$ measures disclosure: $e_\ell$ is the risk of floor and ceiling responses, and a small $\sigma_{\armC}^2$ is compatible with counts that are very informative about $Z_i$. With $\pi=0.2$, the two negatively correlated pairs of \cref{sec:sens-var} cut $L^*$ from about $27$ to $12$. They also double the mutual information between the treatment count and $Z_i$, from $0.11$ to $0.24$ bits, and raise the share of treated respondents whose posterior $\Pr(Z_i=1\mid Y_i,\armT)$ is at least $0.5$ from $11\%$ to $30\%$. At a fixed floor $v=1$, lengthening the list lowers the exposure but leaves the mutual information essentially unchanged ($0.11$ bits for $m=2,\dots,5$). The mutual information also rises as $\rho\to1$, because the parity of the count then becomes informative, so the admissible set $R$ in (d)(iii) is in general an interval. Negatively correlated lists trade disclosure for efficiency. The plan therefore bounds disclosure measures directly, not only the variance of the control count (\cref{sec:design-list}).

\begin{corollary}[Paired versus split]\label{cor:le-paired-split}
Let $c'_{\armL}$ be the cost of a list interview without the direct question, and suppose the control count is independent of $(Z_i,D_i)$, \cref{asm:le-mono} holds and $\pi\le1/2$. Then, for any $c_{\armD}\ge0$,
\[
  \frac{V^*_{\mathrm{paired}}}{V^*_{\mathrm{split}}}\le\frac{c_{\armL}}{c'_{\armL}}\Bigl\{P_0+\frac{(\pi^H)^2P_1}{P_0(\sigma_{\armT}+\sigma_{\armC})^2}\Bigr\}.
\]
\end{corollary}

The paired design therefore wins whenever the direct question adds less than a fraction of about $P_1/P_0$ to the cost of a list interview. In the illustration above, with nearly free direct interviews and equal list costs in the two designs, the ratio of optimal variances is about $0.81$. The cost that the corollary does not price is the risk to \cref{asm:le}(iv). A direct question asked after the list may be answered differently, and a list answered after a direct denial may be pulled towards consistency. This is why the plan of \cref{sec:design-list} keeps the split design as primary and pairs only a subsample at the end of the questionnaire.

\emph{Proof of \cref{thm:design}.} Throughout, the quantities $\nu_0,b_{\armT},b_{\armC},v_a$ are taken to be positive, and $c_{\armD}>0$.

\emph{Part (a).} For any feasible allocation, by the Cauchy--Schwarz inequality,
\[
  K_{\mathrm S}^2=\Bigl(\sum_a\sqrt{v_a/n_a}\,\sqrt{c_an_a}\Bigr)^2\le\Bigl(\sum_av_a/n_a\Bigr)\Bigl(\sum_ac_an_a\Bigr)=B\,\Var\bigl(\widehat{\SDB}_{\mathrm S}\bigr),
\]
with equality if and only if $\sqrt{v_a/n_a}$ is proportional to $\sqrt{c_an_a}$, that is, $n_a\propto\sqrt{v_a/c_a}$. The budget fixes the constant, which gives \eqref{eq:le-neyman}. With equal costs, $n_a\propto\sqrt{v_a}$.

\emph{Part (b).} Use the variables $(n,n_{\armT},n_{\armC})$, with $n_{\armD}=n-n_{\armT}-n_{\armC}$. The problem is to minimise $\Psi=\nu_0/n+b_{\armT}/n_{\armT}+b_{\armC}/n_{\armC}$ subject to $c_{\armD}n+c_+(n_{\armT}+n_{\armC})=B$ and $n\ge n_{\armT}+n_{\armC}$. The objective is strictly convex on the positive orthant and the constraints are linear, so there is a unique minimiser. If $c_+=0$, the budget fixes $n=B/c_{\armD}$ and $\Psi$ decreases in $n_{\armT}$ and $n_{\armC}$, so the constraint $n\ge n_{\armT}+n_{\armC}$ binds. If $c_+>0$, drop that constraint. The relaxed problem has the form of part (a) with costs $(c_{\armD},c_+,c_+)$ attached to $(n,n_{\armT},n_{\armC})$, so its unique minimiser is \eqref{eq:le-paired-int} with value $K_{\mathrm P}^2/B$. This minimiser satisfies $n>n_{\armT}+n_{\armC}$ if and only if $\sqrt{\nu_0/c_{\armD}}>(\sqrt{b_{\armT}}+\sqrt{b_{\armC}})/\sqrt{c_+}$, that is, $\nu_0c_+>c_{\armD}\bar b^2$. In that case it is feasible for the original problem and hence optimal. Otherwise the original minimiser has $n=n_{\armT}+n_{\armC}$. If it did not, the constraint would be slack in a neighbourhood, so the point would be a local and, by convexity, a global minimiser of the relaxed problem; this contradicts uniqueness. On the face $n=n_{\armT}+n_{\armC}=:n_{\armL}$ the budget gives $n_{\armL}=B/c_{\armL}$. By Cauchy--Schwarz, $b_{\armT}/n_{\armT}+b_{\armC}/n_{\armC}\ge\bar b^2/n_{\armL}$, with equality if and only if $n_{\armT}/n_{\armC}=\sqrt{b_{\armT}/b_{\armC}}$. The minimum is therefore $(\nu_0+\bar b^2)/n_{\armL}=c_{\armL}(\nu_0+\bar b^2)/B$.

\emph{Part (c).} Under the floor-and-ceiling model, the treatment-arm response is $Y_i=W_i+Z_i-\ind{Z_i=1,W_i=J}\chi^{\mathrm c}_i-\ind{Z_i=1,W_i=0}\chi^{\mathrm f}_i$, where $\chi^{\mathrm c}_i\sim\Bern(\phi_{\mathrm c})$ and $\chi^{\mathrm f}_i\sim\Bern(\phi_{\mathrm f})$ are independent of everything else. Control and direct responses are unchanged. Hence $\E(\overline Y_{\armT}-\overline Y_{\armC}-\overline D_{\armD})=\SDB-\phi_{\mathrm c}\Pr(W_i=J,Z_i=1)-\phi_{\mathrm f}\Pr(W_i=0,Z_i=1)$, which is \eqref{eq:le-fcbias} for the split design. In the paired design the same computation conditional on $D_i=0$ shifts the probability limit of $\widehat\Delta_0$ by $-\phi_{\mathrm c}\Pr(W_i=J,Z_i=1\mid D_i=0)-\phi_{\mathrm f}\Pr(W_i=0,Z_i=1\mid D_i=0)$. Multiplying by $P_0$ and using $\Pr(Z_i=1,D_i=0)=\pi^H$ gives the stated form.
For the comparative statics, let $\ell$ be optimal at $B$ and $\ell'$ optimal at $B'>B$. Optimality gives
\[
  \frac{V^*_\ell}{B}+\operatorname{bias}_\ell^2\le\frac{V^*_{\ell'}}{B}+\operatorname{bias}_{\ell'}^2,\qquad
  \frac{V^*_{\ell'}}{B'}+\operatorname{bias}_{\ell'}^2\le\frac{V^*_{\ell}}{B'}+\operatorname{bias}_{\ell}^2.
\]
Adding the two inequalities gives $(V^*_\ell-V^*_{\ell'})(1/B-1/B')\le0$, so $V^*_{\ell'}\ge V^*_\ell$. The second inequality then gives $\operatorname{bias}_{\ell'}^2\le\operatorname{bias}_\ell^2+(V^*_\ell-V^*_{\ell'})/B'\le\operatorname{bias}_\ell^2$. If the control count is independent of $(Z_i,D_i)$, the conditional probabilities in \eqref{eq:le-fcbias} equal $\Pr(W_i=J)$ and $\Pr(W_i=0)$, whose sum is $e_\ell$ (the events are disjoint because $J\ge1$), and the bound follows.

\emph{Part (d).} (i) The exposure $p^J+(1-p)^J$ is symmetric about $1/2$, and for $J\ge2$ its second derivative $J(J-1)\{p^{J-2}+(1-p)^{J-2}\}$ is positive, so $p=1/2$ is its unique minimiser. The variance $Jp(1-p)$ is maximised at $p=1/2$.
(ii) For a pair $(X,X')$ with prevalences $1/2$ and correlation $\rho$, $\Pr(X=X'=1)=\E XX'=1/4+\rho/4=:p_{11}$, so $\Pr(X=1,X'=0)=\Pr(X=0,X'=1)=1/2-p_{11}$ and $\Pr(X=X'=0)=p_{11}$. The pair sum takes the values $0,1,2$ with probabilities $p_{11},1-2p_{11},p_{11}$, and has variance $2p_{11}=(1+\rho)/2$. Summing $m$ independent pairs gives \eqref{eq:le-pairs}. Both expressions are strictly increasing in $\rho$. With the control count independent of $(Z_i,D_i)$, the arm variances that enter $V^*_\ell$ are $v_{\armC}=\sigma_{\armC}^2$ and $v_{\armT}=\sigma_{\armC}^2+\pi(1-\pi)$ (split), or $v_{\armC0}=\sigma_{\armC}^2$ and $v_{\armT0}=\sigma_{\armC}^2+\Delta_0(1-\Delta_0)$ (paired). For every fixed allocation the variance is therefore strictly increasing in $\sigma_{\armC}^2$. If $\rho<\rho'$, with optimal allocations $x(\rho)$ and $x(\rho')$, then $V^*(\rho)\le V(x(\rho');\rho)<V(x(\rho');\rho')=V^*(\rho')$. The bias $|\operatorname{bias}_\ell|=(\phi_{\mathrm c}+\phi_{\mathrm f})\,\varpi\,p_{11}^m$, with $\varpi=\pi$ (split) or $\pi^H$ (paired), is non-decreasing in $\rho$. So $\mathrm{MSE}_\ell(B)$ is strictly increasing in $\rho$.
(iii) At $\rho=-1$, $p_{11}=0$, so every pair sums to one and $W_i\equiv m$. By (ii), $\mathrm{MSE}_\ell(B)$ is strictly increasing in $\rho$, so over any closed admissible set $R\subset(-1,1]$ it is minimised at $\min R$. The floor $m(1+\rho)/2\ge v$ gives $R=[2v/m-1,1]$ with $2v/m-1\in(-1,1]$, so $\rho^*=2v/m-1$. There $\sigma_{\armC}^2=v$, and the exposure is $2p_{11}^m$ with $p_{11}=(1+\rho^*)/4=v/(2m)$. Finally, $\frac{\mathrm d}{\mathrm dm}\,m\log\{v/(2m)\}=\log\{v/(2m)\}-1<0$ because $v\le m$, so the exposure decreases in $m$. \qed

\emph{\Cref{cor:le-paired-split}.} The allocation with an empty direct arm is always feasible in the paired design, so $V^*_{\mathrm{paired}}\le c_{\armL}(\nu_0+\bar b^2)$ by part (b). In the split design, part (a) with list cost $c'_{\armL}$ gives $V^*_{\mathrm{split}}=\{\sqrt{v_{\armD}c_{\armD}}+\sqrt{c'_{\armL}}(\sigma_{\armT}+\sigma_{\armC})\}^2\ge c'_{\armL}(\sigma_{\armT}+\sigma_{\armC})^2$. With the control count independent of $(Z_i,D_i)$, $v_{\armC0}=\sigma_{\armC}^2$, $v_{\armT0}=\sigma_{\armC}^2+\Delta_0(1-\Delta_0)$ and $\sigma_{\armT}^2=\sigma_{\armC}^2+\pi(1-\pi)$. Under \cref{asm:le-mono}, $\Delta_0=(\pi-\pi^D)/(1-\pi^D)\le\pi\le1/2$, and $x\mapsto x(1-x)$ is increasing on $[0,1/2]$, so $v_{\armT0}\le\sigma_{\armT}^2$. Hence $\bar b=\sqrt{P_0}(\sqrt{v_{\armT0}}+\sqrt{v_{\armC0}})\le\sqrt{P_0}(\sigma_{\armT}+\sigma_{\armC})$. Dividing, and using $\nu_0=\Delta_0^2P_0P_1=(\pi^H)^2P_1/P_0$, gives the bound. \qed

\subsection{Linking the list experiment to the measurement model}\label{app:pd-irt}

\emph{Proof of \cref{prop:list-irt}.} (a) Since $\delta_j\ge0$, $\{Y^*_{ij}>b_{jc^*}+\delta_j\}\subseteq\{Y^*_{ij}>b_{jc^*}\}$, so $D_{ij}=1$ implies $Z_{ij}=1$, which is \cref{asm:le-mono}. For each $c$, $\Pr(Y_{ij}\ge c\mid\bm\theta_i)=\Pr(\varepsilon_{ij}>b_{jc}+\delta_j-\eta_{ij})=F(\eta_{ij}-b_{jc}-\delta_j)$ by the symmetry of $F$. This is the graded response model with the same loadings and thresholds $b_{jc}+\delta_j$, which remain ordered. The same computation without $\delta_j$ gives $\Pr(Z_{ij}=1\mid\bm\theta_i)=F(\eta_{ij}-b_{jc^*})$. Finally, $\pi^H_j(\bm\theta_i)=\Pr(b_{jc^*}-\eta_{ij}<\varepsilon_{ij}\le b_{jc^*}+\delta_j-\eta_{ij}\mid\bm\theta_i)$, which is \eqref{eq:le-irt-eta}, and $\SDB_j=\Pr(Z_{ij}=1,D_{ij}=0)=\E\,\pi^H_j(\bm\theta_i)$ by \cref{thm:joint-sdb}(a) under \cref{asm:le-mono}.

(b) The list experiment identifies $\pi_j$ (\cref{prop:list-var}). The law of $\eta_{ij}$ is identified from the loadings and the law of $\bm\theta_i$, and $b^D_{jc^*}$ is identified. By (a), $\pi_j=\E F(\eta_{ij}-b_{jc^*})=\psi(\delta_j)$ with $\psi(\delta)=\E F(\eta_{ij}-b^D_{jc^*}+\delta)$. The function $\psi$ is continuous and strictly increasing, because $F$ is, and by dominated convergence it tends to $0$ and $1$ as $\delta\to\mp\infty$. Since $0<\pi_j<1$, the root of $\psi(\delta)=\pi_j$ exists and is unique, and it equals $\delta_j$. The remaining quantities are explicit functions of $\delta_j$ and identified parameters, and $\SDB_j(x)=\E\{\pi^H_j(\bm\theta_i)\mid X_i=x\}$ uses the identified law of $\bm\theta_i$ given $X_i$. A negative root contradicts $\delta_j\ge0$.

(c) If $\delta_j$ is common to all groups, then by (a) the direct item has the same loadings in every group as the private item and thresholds shifted by the same $\delta_j$ in every group, so it is invariant whenever the private item is. $\ECV_j$ and $\ECV^{(g)}_j$ in \eqref{eq:model-ecv} depend only on the loadings and the factor variances, which are therefore the same in the two models. If $\delta_j=\delta_j^{(g)}$ in group $g$, the direct thresholds are $b_{jc}+\delta_j^{(g)}$, so the between-group difference in every threshold is $\delta_j^{(g)}-\delta_j^{(g')}$, with equal loadings: uniform DIF. Applying (b) within group $g$ identifies $\delta_j^{(g)}$ when the group's latent law is identified.
\qed

\begin{proposition}[Overidentification and non-uniform shifts]\label{prop:list-irt-ext}
Under \cref{def:le-shift} and the conditions of \cref{prop:list-irt}(b), the following hold.
\begin{enumerate}[(a),start=4]
  \item \emph{Overidentification.} With $m\ge2$ covariate values, the $m$ equations $\E\{F(\eta_{ij}-b^D_{jc^*}+\delta)\mid X_i=x_l\}=\pi_j(x_l)$ share one unknown, which gives $m-1$ testable restrictions of the uniform shift.
  \item \emph{Non-uniform shift.} Suppose instead that $Z_{ij}=\ind{\upsilon\eta_{ij}+\varepsilon_{ij}\ge b}$ with $\upsilon>0$, that $\varepsilon_{ij}$ is standard normal (the normal-ogive version), and that $\eta_{ij}\mid X_i=x_l\sim\mathrm N(\bar\eta_l,s_\eta^2)$ for $l=1,2$ with $\bar\eta_1\ne\bar\eta_2$. Then $(\upsilon,b)$ is identified by
  \begin{equation}\label{eq:le-irt-nonunif}
    \frac{\upsilon}{\sqrt{1+\upsilon^2s_\eta^2}}=\frac{\Phi^{-1}\{\pi_j(x_1)\}-\Phi^{-1}\{\pi_j(x_2)\}}{\bar\eta_1-\bar\eta_2},\qquad
    b=\upsilon\bar\eta_1-\sqrt{1+\upsilon^2s_\eta^2}\;\Phi^{-1}\{\pi_j(x_1)\},
  \end{equation}
  provided the right-hand side of the first equation lies in $(0,1/s_\eta)$; otherwise the model is rejected.
  \Cref{asm:le-mono} then holds for every $\bm\theta_i$ only if $\upsilon=1$, and the private and direct versions of $\ECV_j$ coincide because the loadings are rescaled jointly.
\end{enumerate}
\end{proposition}

\emph{Proof.} (d) Applying \cref{prop:list-irt}(b) conditionally on $X_i=x_l$ gives a unique root $\delta(x_l)$ for each $l$. The uniform shift implies $\delta(x_1)=\dots=\delta(x_m)$, which is $m-1$ restrictions on identified quantities.

(e) With standard normal $\varepsilon_{ij}$ and $\eta_{ij}\mid X_i=x_l\sim\mathrm N(\bar\eta_l,s_\eta^2)$, $\pi_j(x_l)=\Pr(\varepsilon_{ij}\ge b-\upsilon\eta_{ij})=\Pr(\varepsilon_{ij}+\upsilon\eta_{ij}\ge b)=\Phi\{(\upsilon\bar\eta_l-b)/\sqrt{1+\upsilon^2s_\eta^2}\}$, since $\varepsilon_{ij}+\upsilon\eta_{ij}\sim\mathrm N(\upsilon\bar\eta_l,1+\upsilon^2s_\eta^2)$. Taking $\Phi^{-1}$ and differencing over $l$ gives the first equation of \eqref{eq:le-irt-nonunif}. The map $\upsilon\mapsto\upsilon/\sqrt{1+\upsilon^2s_\eta^2}$ is continuous and strictly increasing from $(0,\infty)$ onto $(0,1/s_\eta)$, since its derivative is $(1+\upsilon^2s_\eta^2)^{-3/2}>0$. So $\upsilon$ is identified when the right-hand side lies in that interval, and $b$ follows from the equation for $l=1$. For monotonicity, $\{\eta+\varepsilon\ge b^D\}\subseteq\{\upsilon\eta+\varepsilon\ge b\}$ for all $(\eta,\varepsilon)$ requires $b^D-\eta\ge b-\upsilon\eta$ for all $\eta\in\mathbb R$. This holds only if $\upsilon=1$ (and then if and only if $b\le b^D$). The private item has loadings $\upsilon(a_{jG},a_{js})$, so its $\ECV_j$ is unchanged. \qed

\subsection{Supplementary results and proofs for \texorpdfstring{\cref{sec:smallarea}}{Section 6}}\label{app:pd-smallarea}

\begin{proposition}[Count indices under non-invariance and nonlinearity]\label{prop:index-dif}
In the setting of \cref{prop:index}, the following hold.
\begin{enumerate}[(a),start=3]
  \item \emph{Uniform DIF.} In a location family with $\mu_1>\mu_2$, invariant slopes $a_j>0$, and $b_j^{(1)}=b_j^{(2)}+\Delta b_j$: if $\Delta b_j=a_j(\mu_1-\mu_2)$ for all $j$, the responses $(U_{i1},\dots,U_{iM})$ have the same law in both groups; if $\Delta b_j\ge a_j(\mu_1-\mu_2)$ for all $j$, with strict inequality for some $j$, then $P_1(q)<P_2(q)$ for every $q$.
  \item \emph{Non-uniform DIF.} Let $M\ge2$, let $F$ and $\Gamma_0$ be symmetric about zero with $\Gamma_0$ continuous, and let $b_j^{(g)}=0$ for all $j$ and $g$. For every $\mu_1>0=\mu_2$ there are slopes $a^{(1)}>0$ (small) and $a^{(2)}>0$ (large), common to all items within a group, such that $P_1(1)>P_2(1)$ but $P_1(M)<P_2(M)$. The sign of the index difference then depends on $q$, although the thresholds are invariant.
  \item \emph{Nonlinearity.} Under full invariance with all $a_j>0$, in a location family, $\mu\mapsto P(q;\mu)$ is strictly increasing from $0$ to $1$. Hence for any $\Delta_A>\Delta_B>0$ there are baselines $\mu_A,\mu_B$ with $P(q;\mu_A+\Delta_A)-P(q;\mu_A)<P(q;\mu_B+\Delta_B)-P(q;\mu_B)$: a larger latent change can produce a smaller change in the index.
\end{enumerate}
\end{proposition}

Part (c) shows that a uniform shift of all thresholds by $a_j(\mu_1-\mu_2)$ is observationally equivalent to equal latent means; this is the index counterpart of \cref{thm:dif-bias}(d), and it is the indeterminacy that anchor items resolve (\cref{prop:anchor}). Slightly larger shifts reverse the sign of every index comparison. Part (d) shows that differences in discrimination alone, with identical thresholds, can make the low-threshold and high-threshold versions of the index disagree about which group has more agreement. Part (e) concerns changes. If two groups start from different baselines, the index can rank their latent changes in the wrong order even under full invariance.

\begin{example}[Numerical instances]\label{ex:sa-index}
Take $M=10$, logistic $F$ and normal $\Gamma_g$. (i) For \cref{prop:index}(b), let all $a_j=1.5$, let the $b_j$ be equally spaced on $[0.5,3.5]$, and let $\Gamma_1=\mathrm N(0.2,0.6^2)$ and $\Gamma_2=\mathrm N(0,1.4^2)$. Then $P_1(q)>P_2(q)$ for $q\le2$ and $P_1(q)<P_2(q)$ for $q\ge3$; at $q=5$, $P_1=0.10$ and $P_2=0.21$. (ii) For part (d) of \cref{prop:index-dif}, let $b_j=0$, $a^{(1)}=0.25$, $a^{(2)}=3$, $\mu_1=0.3$, $\mu_2=0$ and unit variances. Then $P_1(1)=0.999$ and $P_2(1)=0.806$, but $P_1(10)=0.002$ and $P_2(10)=0.194$; the sign changes between $q=5$ and $q=6$. (iii) For part (e), in the limit of perfectly discriminating statements with common location $t$ and $\Gamma_0=\mathrm N(0,1)$, $P(q;\mu)=\Phi(\mu-t)$. A group starting at $\mu=t-3$ whose latent mean rises by $0.5$ moves the index by $0.005$; a group starting at $\mu=t$ whose mean rises by $0.25$ moves it by $0.099$.
\end{example}

\emph{Poststratification identities.} Let a member of area $r$ be drawn at random. Its cell $k$ has probability $w_k$, and given $k$ its latent vector has mean $\bm\mu_k$ and covariance $\bm\Sigma_k$. The law of total expectation gives $\E_r\varphi(\bm\theta)=\sum_kw_k\E[\varphi(\bm\theta)\mid k]$ for every integrable $\varphi$; with $\varphi(\bm\theta)=\bm\theta$ this is \eqref{eq:mrp}, and with $\varphi(\bm\theta)=\ind{\theta_G>t}$ and Gaussian cells it is the tail-share formula. The law of total covariance gives $\bm\Sigma_r=\E_r\{\Var(\bm\theta\mid k)\}+\Var_r\{\E(\bm\theta\mid k)\}$, which is \eqref{eq:mrp-cov}. For a finite population the same identities hold exactly for the empirical distribution of the $N_r$ members, with realised cell means and covariances (divisor $N_k$) in place of $\bm\mu_k$ and $\bm\Sigma_k$. For \cref{rem:sa-finite}: if the members of cell $k$ are independent draws from the superpopulation, the realised cell means are independent with covariance $\bm\Sigma_k/N_k$. So $\sum_kw_k(\bar{\bm\theta}^{\mathrm{pop}}_k-\bm\mu_k)$ has covariance $\sum_kw_k^2\bm\Sigma_k/N_k=N_r^{-1}\sum_kw_k\bm\Sigma_k$, which is of order $N_r^{-1}$.

\emph{\Cref{prop:eap}.} Under the prior $\mathrm N(m_0,\sigma_0^2)$ and $\hat\theta_i\mid\theta_i\sim\mathrm N(\theta_i,s^2)$, the posterior mean is $\tilde\theta_i=m_0+\varrho(\hat\theta_i-m_0)$. Since $\E(\hat\theta_i\mid k)=\E\{\E(\hat\theta_i\mid\theta_i)\mid k\}=\mu_k$, linearity gives \eqref{eq:sa-eap}. The two-step estimator averages $\tilde\theta_i$ within cells, which converges to $\E(\tilde\theta_i\mid k)$, and then poststratifies. Its limit is $\sum_kw_k\{m_0+\varrho(\mu_k-m_0)\}=m_0+\varrho(\bar\theta_r-m_0)$, and the difference between two areas is $\varrho(\bar\theta_r-\bar\theta_{r'})$. In the joint model \eqref{eq:sa-latreg}, $\mu_k$ is a parameter of a correctly specified regular model, and as $n_k\to\infty$ its posterior concentrates at the true value \citep[ch.~10]{vandervaart1998}, so its posterior mean is consistent. In finite cells, partial pooling shrinks it towards the regression prediction by an amount that vanishes as $n_k$ grows, whereas the two-step limit is attenuated by $\varrho$ for every $n_k$.

\emph{\Cref{prop:index}.} Let $V_1,\dots,V_M$ be independent uniform variables, independent of $\theta_i$, and set $U_{ij}=\ind{V_j\le F(a_j\theta_i-b_j)}$. This reproduces \eqref{eq:sa-irf} under full invariance. Write $p_q(\theta)=\Pr(U_{i+}\ge q\mid\theta_i=\theta)$.

\emph{Step 1: $p_q$ is continuous and strictly increasing for $1\le q\le M$ when all $a_j>0$.} Continuity follows from that of $F$. If $\theta<\theta'$, each $U_{ij}$ is pointwise no larger at $\theta$ than at $\theta'$, so $U_{i+}$ is too and $p_q(\theta)\le p_q(\theta')$. Moreover, $0<F<1$ on $\mathbb R$ because $F$ is strictly increasing. So with positive probability exactly $q-1$ of the $V_j$, $j\ne1$, fall below $F(a_j\theta-b_j)$, the others exceed $F(a_j\theta'-b_j)$, and $F(a_1\theta-b_1)<V_1\le F(a_1\theta'-b_1)$. On that event $U_{i+}<q$ at $\theta$ and $U_{i+}\ge q$ at $\theta'$, so $p_q(\theta)<p_q(\theta')$.

\emph{Part (a).} Let $\Gamma_g^{-1}$ be the quantile functions and $\Upsilon$ a uniform variable. Stochastic dominance is equivalent to $\Gamma_1^{-1}(u)\ge \Gamma_2^{-1}(u)$ for all $u\in(0,1)$, and $P_g(q)=\E\,p_q\{\Gamma_g^{-1}(\Upsilon)\}$. Step 1 gives $P_1(q)\ge P_2(q)$. If $\Gamma_1\ne\Gamma_2$, the two quantile functions differ at some $u_0$. Both are left-continuous and non-decreasing, so they differ on a set of positive Lebesgue measure; otherwise they would agree on a dense set and hence everywhere. On that set $p_q\{\Gamma_1^{-1}(u)\}>p_q\{\Gamma_2^{-1}(u)\}$, which makes the inequality strict. In a location family $\Gamma_1^{-1}(u)-\Gamma_2^{-1}(u)=\mu_1-\mu_2$ for all $u$, which gives the sign statement.

\emph{Part (b).} Set $a_j=a$ and $b_j=at$ for all $j$. For $\theta>t$, $p_q(\theta)\ge\Pr(U_{i+}=M\mid\theta)=F\{a(\theta-t)\}^M\to1$ as $a\to\infty$. For $\theta<t$, $p_q(\theta)\le\Pr(U_{i+}\ge1\mid\theta)\le MF\{a(\theta-t)\}\to0$. Since $\Pr(\theta_i=t)=0$, dominated convergence gives $P_g(q)\to1-\Phi\{(t-\mu_g)/\sigma_g\}$ for every $q\in\{1,\dots,M\}$. Now $(t-\mu_1)/\sigma_1>(t-\mu_2)/\sigma_2$ if and only if $t(\sigma_2-\sigma_1)>\mu_1\sigma_2-\mu_2\sigma_1$, that is, $t>t_\times$. So the limit satisfies $P_1(q)<P_2(q)$ for $t>t_\times$ and $P_1(q)>P_2(q)$ for $t<t_\times$, and both strict inequalities persist for all sufficiently large finite $a$. The largeness condition matters: for moderate $a$ the sign at a given $t$ can depend on $q$.

\emph{\Cref{prop:index-dif}, part (c).} Write $\theta_i=\mu_g+\xi_i$, with $\xi_i\sim \Gamma_0$ in both groups. In group 1 the linear predictor of item $j$ is $a_j(\mu_1+\xi_i)-b^{(2)}_j-\Delta b_j=\{a_j(\mu_2+\xi_i)-b^{(2)}_j\}+\{a_j(\mu_1-\mu_2)-\Delta b_j\}$. If $\Delta b_j=a_j(\mu_1-\mu_2)$ for all $j$, the conditional law of $(U_{i1},\dots,U_{iM})$ given $\xi_i$ is the same in both groups, and so is the law of $\xi_i$. If $\Delta b_j\ge a_j(\mu_1-\mu_2)$ for all $j$, each item probability in group 1 is no larger, given $\xi_i$, than in group 2, and it is strictly smaller for some $j^*$. Now $\Pr(U_{i+}\ge q)=p_{j}\Pr(U_{i,-j}\ge q-1)+(1-p_j)\Pr(U_{i,-j}\ge q)$ is affine in each item probability $p_j$ with slope $\Pr(U_{i,-j}=q-1)\ge0$, where $U_{i,-j}$ omits item $j$. That slope is positive when all other item probabilities lie in $(0,1)$. Hence $\Pr(U_{i+}\ge q\mid\xi_i,g=1)<\Pr(U_{i+}\ge q\mid\xi_i,g=2)$ for every $\xi_i$, and integrating over $\Gamma_0$ gives $P_1(q)<P_2(q)$.

\emph{Part (d).} With $b_j=0$ and common slope $a^{(g)}$ in group $g$, $P_g(q)$ is continuous in $a^{(g)}\in[0,\infty)$ by dominated convergence. At $a^{(1)}=0$, the item probabilities equal $F(0)=1/2$ by symmetry, so $U_{i+}\mid g=1\sim\Bin(M,1/2)$, $P_1(1)=1-2^{-M}$ and $P_1(M)=2^{-M}$. As $a^{(2)}\to\infty$, the argument of part (b) with $t=0$ gives $P_2(q)\to\Pr(\theta_i>0\mid g=2)=1/2$ for every $q$, because $\Gamma_0$ is continuous and symmetric and $\mu_2=0$. For $M\ge2$, $1-2^{-M}>1/2>2^{-M}$. By continuity, $P_1(1)>P_2(1)$ and $P_1(M)<P_2(M)$ for all sufficiently small $a^{(1)}>0$ and sufficiently large $a^{(2)}$.

\emph{Part (e).} $P(q;\mu)=\E\,p_q(\mu+\xi_i)$ is strictly increasing in $\mu$ by Step 1. As $\mu\to-\infty$, $p_q(\mu+\xi_i)\le\Pr(U_{i+}\ge1\mid\mu+\xi_i)\le\sum_jF\{a_j(\mu+\xi_i)-b_j\}\to0$ pointwise, and as $\mu\to\infty$, $p_q(\mu+\xi_i)\ge\prod_jF\{a_j(\mu+\xi_i)-b_j\}\to1$. Dominated convergence gives the limits $0$ and $1$. Fix $\mu_B$ and let $c=P(q;\mu_B+\Delta_B)-P(q;\mu_B)>0$. Choose $\mu_A$ with $P(q;\mu_A+\Delta_A)<c$. Then $P(q;\mu_A+\Delta_A)-P(q;\mu_A)<c$.

\emph{Several factors.} If $U_{ij}$ depends on $\bm\theta_i\in\mathbb R^3$ through $F(\bm a_j^\top\bm\theta_i-b_j)$ with $\bm a_j\ge\bm0$ componentwise, the coupling of Step 1 shows that $p_q$ is non-decreasing in each coordinate. The usual multivariate stochastic order is defined by $\E f(\bm\theta^{(1)})\ge\E f(\bm\theta^{(2)})$ for all bounded non-decreasing $f$, so $P_1(q)\ge P_2(q)$ follows directly. \qed
